%==============================================================================
% MUWS 2026 draft #1 — Label-prior collapse / responsible evaluation
% Compile on Overleaf: upload this file + refs.bib, set compiler to pdfLaTeX,
% run twice (+ BibTeX). Uses acmart [sigconf]; switch off 'nonacm' for camera-ready.
%==============================================================================
\documentclass[sigconf,nonacm]{acmart}
\settopmatter{printacmref=false}
\usepackage{booktabs}
\usepackage{graphicx}

\begin{document}

\title{Do Small Vision--Language Models Actually Classify?\\
Label-Prior Collapse in Zero-Shot Multimodal Hate and Sarcasm Detection}

\author{Anonymous Author(s)}
\affiliation{\institution{Anonymous Institution}\country{}}

\begin{abstract}
Small vision--language models (VLMs, $\leq$2B parameters) are increasingly proposed for
on-device moderation of hateful and sarcastic multimodal content. We show that, used
zero-shot, such models largely do \emph{not} classify: they collapse to a near-constant
answer regardless of input. On Hateful Memes and MMSD2.0, SmolVLM-256M/500M answer
``yes'' for 86--98\% of examples while Gemma-4-E2B answers ``no'' for 100\%. Critically,
the field's standard thresholded metrics (accuracy, F1) \emph{mask} this failure---on a
class-skewed slice we observe an apparent F1 of 0.89 that collapses to chance under
balanced evaluation. We argue that ranking-based metrics (AUROC) together with a simple
prediction-distribution diagnostic are necessary to report small-VLM behaviour honestly,
and we show that the principal effect of light LoRA fine-tuning is \emph{de-biasing}:
prediction-positive rates move from $\sim$0.95 to $\sim$0.30 (close to the true base
rate), recovering genuine discrimination. We release a reproducible evaluation harness
and recommend reporting practices for multimodal human-understanding tasks.
\end{abstract}

\keywords{multimodal hate speech, sarcasm detection, vision--language models,
evaluation, calibration, bias}

\maketitle

\section{Introduction}
Detecting hateful and sarcastic content in image--text memes is a central task for
multimodal human understanding on the social web. As compact VLMs become attractive for
privacy-preserving, on-device moderation, a natural question is how well small models
($\leq$2B parameters) perform these tasks out of the box. We report a cautionary
empirical finding: in the zero-shot regime, small VLMs frequently do not behave as
classifiers at all. Instead they exhibit \emph{label-prior collapse}---answering the
same label for nearly every input---and this collapse is hidden by the thresholded
metrics commonly reported on these benchmarks.

Our contributions are: (i) a systematic demonstration of label-prior collapse across
three small VLMs and two tasks, including a worked example where a class-skewed
evaluation slice produces a misleadingly high F1; (ii) evidence that AUROC plus a
prediction-distribution diagnostic is the minimal honest reporting standard for these
models; and (iii) an analysis showing that the dominant effect of light parameter-efficient
fine-tuning is calibration/de-biasing rather than a large gain in separability for the
harder task. All experiments use an open, resumable harness.

\section{Related Work}
The Hateful Memes Challenge~\cite{kiela2020hateful} was explicitly designed with
``benign confounders'' so that unimodal or weak models fail, making it a stress test for
multimodal reasoning. MMSD2.0~\cite{qin2023mmsd2} re-annotates the multimodal sarcasm
corpus of \citet{cai2019multimodal} to remove spurious textual cues. Small instruction-tuned
VLMs such as SmolVLM~\cite{marafioti2025smolvlm}, built on Idefics3~\cite{laurencon2024idefics3},
target efficient deployment. Calibration of neural classifiers is well studied in the
unimodal setting~\cite{guo2017calibration}; we connect that lens to zero-shot multimodal
moderation and show that prediction-distribution collapse, not just miscalibrated
confidence, is the first-order problem.

\section{Setup}
\textbf{Models.} SmolVLM-256M-Instruct, SmolVLM-500M-Instruct, and Gemma-4-E2B
($\sim$2B effective parameters). \textbf{Tasks.} Hateful Memes (hateful vs.\ not; dev
split, balanced, 500 examples) and MMSD2.0 (\texttt{mmsd-v2}; sarcastic vs.\ not; test
split, 500 examples). \textbf{Protocol.} We prompt each model with a task instruction,
the image, and a yes/no question, and read $P(\text{yes})$ from the softmax over the
yes/no answer tokens---a continuous score for AUROC---thresholding at 0.5 for hard
labels. We report AUROC (rank-based, threshold-free), accuracy, and the
\emph{prediction-positive rate} (\texttt{pred\_pos}, the fraction of examples assigned
the positive label). All runs use a fixed seed and a single NVIDIA A100; the harness
loads each model once and appends one record per cell to a JSONL leaderboard.

\section{Results}

\subsection{Zero-shot small VLMs collapse to a constant answer}
Table~\ref{tab:zs} shows that no zero-shot model meaningfully separates the classes:
AUROC ranges from 0.50 to 0.60. The \texttt{pred\_pos} column reveals the mechanism:
the SmolVLM models answer ``yes'' for 86--98\% of inputs, whereas Gemma-4-E2B answers
``no'' for \emph{every} input (and, despite $\sim$4$\times$ the parameters and $\sim$7$\times$
the memory, does not separate the classes better). Accuracy is therefore at or below the
majority-class baseline and is uninformative.

\begin{table}[t]
\centering
\caption{Zero-shot (fp16, $n{=}500$). AUROC is the only metric that distinguishes faint
signal from collapse; \texttt{pred\_pos} exposes the constant-answer behaviour.}
\label{tab:zs}
\begin{tabular}{llccc}
\toprule
Dataset & Model & AUROC & Acc. & \texttt{pred\_pos} \\
\midrule
Hateful Memes & SmolVLM-256M & 0.513 & 0.506 & 0.98 \\
              & SmolVLM-500M & 0.581 & 0.532 & 0.93 \\
              & Gemma-4-E2B  & 0.528 & 0.500 & 0.00 \\
\midrule
MMSD2.0       & SmolVLM-256M & 0.502 & 0.414 & 0.86 \\
              & SmolVLM-500M & 0.600 & 0.424 & 0.98 \\
              & Gemma-4-E2B  & 0.527 & 0.584 & 0.00 \\
\bottomrule
\end{tabular}
\end{table}

\subsection{Thresholded metrics hide the collapse}
Because the models predict one label almost always, thresholded scores are dominated by
the class balance of the evaluation slice. On a small, positive-skewed slice
($n{=}50$) of MMSD2.0 we measured an apparent \textbf{F1 of 0.89} for SmolVLM-256M; on a
balanced $n{=}500$ slice the same configuration yields AUROC 0.50 and accuracy below the
majority baseline. The F1 was an artifact of a model that says ``yes'' meeting a slice
that is mostly positive. This is the core reporting hazard: \emph{F1/accuracy on a
collapsed predictor measure the test set, not the model.}
% TODO Fig.1: prediction-score histograms per model (mass piled at 0 or 1) +
% reliability diagram. Generate with scripts/make_figures style code over the JSONL.
\begin{figure}[t]
\centering
% \includegraphics[width=\columnwidth]{figs/collapse_histograms.pdf}
\fbox{\parbox{0.9\columnwidth}{\centering\vspace{1.2cm}\textit{Figure 1
(placeholder): per-model $P(\text{yes})$ histograms and reliability diagram showing
score mass collapsed near 0 or 1.}\vspace{1.2cm}}}
\caption{Prediction-distribution diagnostic. Score mass piles at one extreme,
not spread across the decision boundary.}
\label{fig:hist}
\end{figure}

\subsection{LoRA's first-order effect is de-biasing}
We apply LoRA~\cite{hu2022lora} ($r{=}16$, attention+MLP projections, 1.85\% of
parameters) to SmolVLM-500M, fine-tuning it to emit yes/no on 3{,}000 in-domain examples
(2 epochs). Table~\ref{tab:lora} shows that fine-tuning moves \texttt{pred\_pos} from
$\sim$0.95 toward the true base rate ($\sim$0.42--0.50) and that the model begins to
discriminate. The effect is dramatic for sarcasm (AUROC $0.600\to0.881$) and modest for
the adversarially hard Hateful Memes ($0.581\to0.619$)---but in both cases the
\emph{qualitative} change is the same: the model stops answering with a constant label.

\begin{table}[t]
\centering
\caption{In-domain LoRA on SmolVLM-500M (fp16, $n{=}500$). Fine-tuning de-biases
(\texttt{pred\_pos}$\to$base rate) before it improves separability.}
\label{tab:lora}
\begin{tabular}{llcc}
\toprule
Dataset & Condition & AUROC & \texttt{pred\_pos} \\
\midrule
Hateful Memes & zero-shot & 0.581 & 0.93 \\
              & + LoRA    & 0.619 & 0.36 \\
\midrule
MMSD2.0       & zero-shot & 0.600 & 0.98 \\
              & + LoRA    & 0.881 & 0.29 \\
\bottomrule
\end{tabular}
\end{table}

\section{Discussion and Recommendations}
Our results suggest concrete reporting practices for small-VLM multimodal moderation:
(1) always report AUROC alongside, or instead of, thresholded F1/accuracy; (2) report
the prediction-positive rate (or full score histogram) so collapse is visible; (3)
evaluate on class-balanced or full splits and avoid small skewed slices; and (4) when
claiming that fine-tuning ``improves'' a small VLM, separate the de-biasing component
from genuine gains in separability. These practices are cheap and would prevent the kind
of inflated headline number we document.

\section{Limitations}
We study two datasets, three models, and single-seed runs on $n{=}500$ balanced eval
slices; confidence intervals are approximately $\pm0.03$--$0.04$ AUROC. We do not vary
the prompt extensively; Gemma-4-E2B in particular ships without a chat template and is
evaluated with a raw prompt, so its collapse may partly reflect prompt sensitivity rather
than capability. Our claims concern the zero-shot small-VLM regime and light LoRA, not
large or fully fine-tuned systems.

\section{Conclusion}
Small zero-shot VLMs for multimodal hate and sarcasm detection tend to collapse to a
constant answer, and standard thresholded metrics conceal this. Reporting AUROC with a
prediction-distribution diagnostic exposes the behaviour, and reframes the value of
light fine-tuning as primarily one of de-biasing. We hope these practices improve how
small multimodal moderation models are evaluated.

\bibliographystyle{ACM-Reference-Format}
\bibliography{refs}
\end{document}
